\documentclass[letter,11pt]{article}
\usepackage{hyperref}
\hypersetup{pdfstartview={XYZ null null 1.00}}
\usepackage{graphicx}
\usepackage{fancyhdr}
\usepackage{float} % Force figure placement
\usepackage{multicol}
\pagestyle{fancy}
\usepackage[letterpaper, margin=1in]{geometry}
\geometry{letterpaper}
\usepackage{parskip} % Disable initial indent
\usepackage{attachfile}

\usepackage[utf8]{inputenc}
\fancyhf{}
\renewcommand{\headrulewidth}{0pt} % Remove default underline from header package
\rhead{CYBR/CMSC 691 - Final Lab}
%\rhead{}
\lhead{\begin{picture}(0,0) \put(0,-10){\includegraphics[width=1.1cm]{Images/UMBC-vertical}} \end{picture}}
\cfoot{\thepage}
\rfoot{Spring 2020}
\lfoot{Final Lab}
%\AtEndDocument{\vfill \hfill \LaTeX}

\begin{document}

\section*{Final Lab: Feature Engineering}
\paragraph{}Good science, and the scientific method, are about reproducability in experiments. In their paper, \textit{Feature Selection for Malware Detection Based on Reinforcement Learning} \attachfile{Papers/Malware_Feature_Selection_Reinforcement_Learning.pdf}, Fang, Z., et al discuss domain-specific features and specific machine learning algorithms used, along with achieved accuracy results. Using the same features, and the malware samples we've been using this semester, try to reproduce their results from their paper. The results are in the table on Page 9.

\paragraph{}Using any prior lab as a starting point, create new datasets which are based on the features discussed in the paper. Refer to documentation for the various algorithms in Sklearn:
\begin{itemize}
    \item Decision Tree: \url{https://scikit-learn.org/stable/modules/tree.html}
    \item Random Forest: \url{https://scikit-learn.org/stable/modules/generated/sklearn.ensemble.RandomForestClassifier.html}
    \item Naive Bayes: \url{https://scikit-learn.org/stable/modules/naive_bayes.html}
    \item SVM: \url{https://scikit-learn.org/stable/modules/generated/sklearn.svm.SVC.html}
\end{itemize}

\paragraph{}You may skip the KNN algorithm, since we didn't discuss it. All algorithms use the same function names. Generally they follow this format:
\begin{itemize}
    \item \texttt{model = ModelName()}
    \item \texttt{model.fit(data, dataLabels)}
    \item \texttt{accuracy = model.score(testData, testLabels)}
\end{itemize}

\paragraph{}Please email me with any questions.

\textbf{Points:} 50\\
\textbf{Reminder:} This is part of the final exam.

\section*{Questions}
\paragraph{}In a separate document, provide your answers to the following questions. The expected format is a Word document, or similar, with a few images demonstrating your progress. The assignment will be submitted via Blackboard.

\begin{itemize}
    \item Which machine learning algorithm was most effective?
    \item Where they any difficulties you encountered recreating their results?
    \item Do you think the paper was accurate in their claims (were you able to reproduce their accuracy results)?
\end{itemize}
\end{document}